% SPDX-License-Identifier: Apache-2.0
%
% Vreteno through an open ASIC flow. Built by //docs:tapeout. Every
% listing is included from the tree and every number comes from a run.
\documentclass[journal,9pt]{IEEEtran}

\newcommand{\listingsize}{\fontsize{6}{7}\selectfont}
% SPDX-License-Identifier: Apache-2.0
%
% The house preamble, shared by every document in this directory. Loaded
% after \documentclass, before \title. Keeping it in one file is what
% stops two articles drifting apart in font, listing style or figure
% style.
% Computer Modern. IEEEtran selects Times (ptm) by default, so the face has
% to be chosen explicitly.
%
% Latin Modern is the Computer Modern variety used here, and the choice is
% forced rather than aesthetic. Under T1 encoding, plain `cmr` has no Type 1
% outlines unless cm-super is installed, and it is not in the pinned
% distribution, so pdflatex falls back to EC bitmaps and every glyph in the
% PDF becomes a Type 3 font. That was measured: the first build produced a
% document with 10 Type 3 fonts and no Type 1. Latin Modern is Computer
% Modern redrawn with a full T1 repertoire and real .pfb outlines, so the
% same design comes out scalable.
\usepackage[T1]{fontenc}
\usepackage[utf8]{inputenc}
\usepackage{lmodern}
% microtype is not in the pinned TeX distribution. emergencystretch alone
% keeps the two-column measure from producing overfull lines.
\setlength{\emergencystretch}{3em}

\usepackage{amsmath}
\usepackage{amssymb}
\usepackage{amsthm}
\usepackage{booktabs}
\usepackage{array}
% enumitem is not in the pinned TeX distribution; the list spacing below
% does what \setlist would have done.
\usepackage{float}
\usepackage{xcolor}
\usepackage{listings}
\usepackage{tikz}
\usetikzlibrary{arrows.meta,positioning,fit,backgrounds,calc}
\usepackage[hidelinks,bookmarks=true]{hyperref}

% Numbered, definition-style box for a rule the reader is meant to apply.
\theoremstyle{definition}
\newtheorem{rulebox}{Rule}
\newtheorem{example}{Example}

% Unnumbered asides.
\theoremstyle{remark}
\newtheorem*{tip}{Tip}
\newtheorem*{pitfall}{Pitfall}

% Tighter lists, in place of enumitem's \setlist.
\setlength{\topsep}{2pt}
\setlength{\itemsep}{1pt}
\setlength{\parsep}{0pt}

% Code in the running text. The typewriter face under T1 ligates `<<`
% and `>>` into guillemets, so a shift printed as \code{<<} came out as
% a French quotation mark (issue 571). microtype's \DisableLigatures is
% not in the pinned distribution, but the primitive it rests on is
% pdfTeX's own: \pdfnoligatures strips every ligature from the font it
% is given, and the current font inside \texttt is the typewriter face
% at the size in use, so each size is stripped the first time \code
% sets it. Code wants no ligature at all: two quotes are two quotes.
\newcommand{\code}[1]{\texttt{\ifdefined\pdfnoligatures\pdfnoligatures\font\fi #1}}
\newcommand{\term}[1]{\emph{#1}}

% Syntax highlighting. Four roles, distinguishable in colour and still
% distinguishable in greyscale, because an IEEE article gets printed: the
% keyword is the only bold role, the comment is the only italic one, and the
% two remaining roles differ in lightness as well as in hue.
\definecolor{kwblue}{RGB}{20,60,150}
\definecolor{cmtgreen}{RGB}{90,120,90}
\definecolor{strred}{RGB}{150,50,40}
\definecolor{numgray}{gray}{0.45}
\definecolor{framegray}{gray}{0.70}
\definecolor{bgshade}{gray}{0.97}
% A darker fill for figure cells. The listing background has to stay very
% light to keep code readable; a figure cell has to be clearly shaded or the
% distinction it draws is invisible in print.
\definecolor{cellshade}{gray}{0.90}

% The listing size is the document's choice, because it depends on the
% column: a two-column article sets `\listingsize` to 6pt and keeps its
% listings within 70 columns, or to 5.6pt when it includes source that
% rustfmt sets at 80, which is what fits a column's frame (issue 1061);
% a one-column document keeps the body size and includes source files
% at 80. `//docs:frame_test` checks every listing line against its frame. Lines are never broken; a line that does not fit
% overflows the frame, which is visible, rather than wrapping, which is
% not.
\providecommand{\listingsize}{\scriptsize}

% `'` is deliberately not a string delimiter: the language uses it for loop
% labels, as Rust does, so treating it as a quote would swallow the rest of
% the line.
\lstdefinestyle{txhdl}{
  basicstyle=\ttfamily\listingsize,
  keywordstyle=\bfseries\color{kwblue},
  commentstyle=\itshape\color{cmtgreen},
  stringstyle=\color{strred},
  numberstyle=\tiny\color{numgray},
  morekeywords={mod,use,pub,const,type,struct,enum,trait,impl,fn,async,let,mut,
    match,if,else,loop,while,for,in,break,continue,return,as,await,
    module,bus,channel,transaction,protocol,pipeline,flow,seq,config,
    port,stage,pipe,instance,connect,bind,tag,domain,blackbox,foreign,
    property,role,signal,handler,true,false,
    sequence,interface,modport,implements,package,import,var,wire,func,proc,
    case,default,next,fork,branch,join,state,fsm},
  morecomment=[l]{//},
  morestring=[b]",
  showstringspaces=false,
  columns=fullflexible,
  keepspaces=true,
  breaklines=false,
  % Framed, so a listing is bounded rather than floating in the column.
  frame=single,
  rulecolor=\color{framegray},
  framesep=3pt,
  backgroundcolor=\color{bgshade},
  xleftmargin=3pt,
  xrightmargin=3pt,
  aboveskip=6pt,
  belowskip=6pt,
  % Every listing is numbered and captioned, so it can be referred to by
  % number. `caption` without `float` numbers the listing and leaves it
  % where it was written, which is what a listing in running prose needs.
  captionpos=b,
  abovecaptionskip=2pt,
  belowcaptionskip=6pt,
}

% Figures drawn as text. Framed the same way, and with the highlighting
% switched off, because a box-and-arrow diagram is not code and half its
% words turning blue reads as an accident. Setting a style adds keys rather
% than clearing the ones already set, so the three roles are neutralised by
% hand rather than by leaving them out. Program output is printed in this
% style too, and a netlist's assignment can run past the frame, so a line
% that does not fit is broken and the rest indented; source never is.
\lstdefinestyle{ascii}{
  basicstyle=\ttfamily\listingsize,
  keywordstyle=\ttfamily,
  commentstyle=\ttfamily,
  stringstyle=\ttfamily,
  columns=fullflexible,
  keepspaces=true,
  breaklines=true,
  breakindent=2em,
  frame=single,
  rulecolor=\color{framegray},
  framesep=3pt,
  backgroundcolor=\color{bgshade},
  xleftmargin=3pt,
  xrightmargin=3pt,
  aboveskip=6pt,
  belowskip=6pt,
}
\lstset{style=txhdl}

% House TikZ styles. `shorten` lives on the arrow styles rather than on each
% arrow, so every arrow in the document gets the gap, including ones added
% later. TikZ clips a path at a node's border, so without it an arrowhead
% butts against the box with no gap at all. Keep `node distance` well above
% twice the shorten, or the arrow shrinks to nothing.
\tikzset{
  box/.style={draw,rounded corners=2pt,align=center,inner sep=4pt,
              font=\footnotesize},
  boxf/.style={box,fill=bgshade},
  lbl/.style={font=\scriptsize\itshape,align=center},
  fl/.style={-{Stealth[length=4pt]},shorten >=2.5pt,shorten <=2.5pt},
  fld/.style={fl,dashed},
}


% The contents list: the class gives a section's number 2.75em, which
% a Roman numeral past thirty overruns onto the title, so the box is
% wider here, and a subsection's entry starts where a title does.
\makeatletter
\def\l@section#1#2{\addpenalty{\@secpenalty}\addvspace{1.0em plus 1pt}%
    \@tempdima 5.75em \begingroup \parindent \z@ \rightskip \@pnumwidth%
    \parfillskip-\@pnumwidth {\bfseries\leavevmode #1}\hfil\hbox to\@pnumwidth{\hss #2}\par%
    \endgroup}
\def\l@subsection{\@dottedtocline{2}{5.75em}{5em}}
\def\l@subsubsection{\@dottedtocline{3}{10.75em}{5.5em}}
\makeatother


\InputIfFileExists{buildstamp.tex}{}{}
\providecommand{\buildstamp}{Draft build (outside the Bazel flow).}

% The numbers the prose quotes, written by the run that produced them.
\input{asic_counts.tex}

% A script of the flow, highlighted for its own vocabulary: the
% comment character is the hash, and the commands are the steps.
\lstdefinestyle{script}{%
  style=txhdl,
  morecomment=[l]{\#},
  morekeywords={set,expr,proc,foreach,source,puts,format,open,close,
    read_verilog,hierarchy,synth,opt,dfflibmap,abc,setundef,splitnets,
    opt_clean,hilomap,insbuf,tee,stat,write_verilog,read_lef,
    read_liberty,read_sdc,link_design,initialize_floorplan,place_pins,
    tapcell,pdngen,global_placement,detailed_placement,check_placement,
    clock_tree_synthesis,set_propagated_clock,repair_design,
    repair_clock_nets,repair_timing,global_route,detailed_route,
    filler_placement,estimate_parasitics,write_def,create_clock,
    set_input_delay,set_output_delay,set_driving_cell,set_load,
    set_max_fanout,set_clock_uncertainty,set_thread_count,get_ports,
    get_clocks,all_inputs,all_outputs,current_design,lsearch,metric,
    report_design_area,report_worst_slack,report_tns,report_checks,
    report_power,report_check_types,report_clock_skew},
}

\title{A Tapeout Prep in TxHDL: Vreteno on an Open 45\,nm Node}

\author{Filip Filmar and Dragiša Janković%
\thanks{Both authors are independent. Filip Filmar is the
corresponding author, at f@hdlfactory.com; Dragiša Janković is
reached at linkedin.com/in/dragisa-jankovic.}%
\thanks{This document examines the trajectory of a design between
synthesis and photomask generation. The target processor is the Vreteno
core described in~\cite{vreteno}, instantiated from the lowered Verilog
emitted by TxHDL. Every listing is included verbatim from the project
repository \code{HDL/txhdl} at build time; the companion
index~\cite{cover} catalogs the documentation suite.}%
\thanks{The authors used a large language model, Claude, as an
assistant in exploring the concepts and in writing and constructing
the documents and the programs; every commit of the repository says
so and carries the prompts that led to it, verbatim.}%
\thanks{\buildstamp}}

\markboth{A Tapeout Prep in TxHDL}{Filmar and Janković: A Tapeout Prep in TxHDL}

\begin{document}
\maketitle

\begin{abstract}
A design that has been simulated, lowered, verified against its execution
trace, and synthesized onto an FPGA has still never taken physical shape
on silicon. This paper takes the Vreteno core's structural Verilog—the
exact netlist targeted to hardware in its FPGA build—maps it onto an open
45\,nm standard-cell library using Yosys, and floorplans, places, routes,
and times it with OpenROAD. The result is a die \asicdie\,$\mu$m square
holding \asiccells{} cells that meets a \asicperiod\,ns clock with
\asicslack\,ns of setup slack and routes with zero DRC violations.
Both toolchains are pinned by cryptographic checksum and unpacked
directly into the build sandbox, requiring no pre-installed EDA packages
on the host. This preparation stops short of a tapeout; the closing
section catalogs what physical fabrication would still demand.
\end{abstract}

\section{What This Is}
\label{sec:t-what}

An FPGA answers a pragmatic question about a design: does it fit, and
how fast does it run, on a specific commercial chip. An ASIC flow
demands things an FPGA never asks for. There are no prefabricated block
RAMs for memories and no dedicated DSP blocks for arithmetic; both must
be synthesized out of raw gates. There is no configuration bitstream to
initialize registers at power-on; a register that was initialized in
elaboration but lacks an explicit hardware reset powers up in an
undefined state. And there is no pre-existing routing fabric: every wire
must be drawn in metal, and its physical length is an unyielding delay.

The Vreteno core has already cleared the FPGA bar~\cite{vreteno}. It runs
on an Alinx AX7A200B development board at 100\,MHz, and its observed
behavior on hardware matches its cycle-accurate simulation bit for
bit. This paper puts that same netlist (as tagged at commit
\code{2c812d3} on 17 September 2026) through an open physical ASIC flow,
reporting the gate count, floorplan, and timing closure achieved.

None of this constitutes a tapeout. No commercial fab operates a process
line corresponding to this library, no photomask was manufactured, and
the rigorous signoff decks demanded by a foundry are absent here.
Section~\ref{sec:t-not} catalogs what remains missing, one requirement
at a time: presenting a layout plot without declaring the verification
gap would claim more than has been achieved.

\section{The Node}
\label{sec:t-node}

Nangate45 is an open standard-cell library targeting a notional 45\,nm
technology node. It provides Liberty (\code{.lib}) timing models for
each cell, LEF abstractions defining cell bounding boxes and pin
geometries, and a technology LEF specifying the design rules, routing
pitches, and wire widths across ten metal layers. It is the reference
library distributed alongside OpenROAD under Apache~2.0. The TxHDL build
fetches the six library files by cryptographic checksum at a pinned
revision, preserving license compliance alongside reproducibility.

It is not a foundry PDK. No physical wafer was ever fabricated from it,
its timing figures reflect plausible 45\,nm predictive models rather
than silicon measurements, and it includes no design rule deck that a fab
would sign off on. Yet those attributes make it the ideal vehicle for an
open reference flow: every algorithmic step is genuine, every reported
metric has physical meaning, and no non-disclosure agreement prevents
publishing the inputs, scripts, and results in a public repository.

What the node lacks is a memory compiler. A commercial foundry kit ships
with software generators that produce dense SRAM macros of whatever width
and depth are required. Nangate45 provides only empty mock macros of
fixed dimensions for placement experimentation, and an empty macro holds
no data. Section~\ref{sec:t-cells} details what the core's memories
became instead—yielding the most instructive architectural lesson of
this run.

\section{The Tools, and How They Are Pinned}
\label{sec:t-tools}

Every toolchain dependency in this repository is fetched by cryptographic
checksum and unpacked into the build sandbox; nothing is assumed to be
installed on the host system. The ASIC toolchain adheres strictly to this
discipline. Achieving hermeticity is uniquely demanding here because
Yosys and OpenROAD are distributed as Debian packages, which declare
shared library dependencies dynamically rather than bundling them.

The build resolves the complete dependency closure in advance.
\code{//tools/debclosure} inspects an upstream distribution's package
index, walks the dependency tree starting from the required binaries,
and writes an immutable lockfile. Each entry records the package name,
exact version, payload size, SHA-256 digest, and dual retrieval URLs.
The first URL points to the live Debian mirror—fast, but ephemeral across
distribution updates; the second points to \code{snapshot.debian.org}—slow,
but permanently immutable. A fetcher that tries the live mirror before
falling back to the snapshot achieves maximum download speed while
guaranteeing bit-for-bit reproducibility years into the future.

Two such lockfiles are maintained in the tree. Yosys 0.52 and its ten
shared-library dependencies comprise eleven packages totaling
13\,MB. OpenROAD—pinned to the final binary release published by the
project—spans 52\,MB in isolation, but its transitive closure demands
131 auxiliary packages totaling an additional 75\,MB. Roughly two-thirds
of that footprint comprises Qt and the associated Mesa, LLVM, X11, and
ICU runtime libraries that Qt pulls in, which the executable dynamically
links regardless of whether an interactive GUI is invoked. Alongside
OpenROAD sits the external linear-programming solver utilized by its
analytical placer, which the upstream package expects in its runtime
search path without bundling.

All package extraction is performed natively by Bazel. Because a
\code{.deb} package is an \code{ar} archive containing a \code{tar.xz}
payload, Bazel's internal decompressors extract both layers in sequence.
The repository rule in \code{//third\_party/debs} arranges the unpacked
artifacts into a unified directory tree against which the EDA tools
execute, isolated with their private shared libraries. The entire flow
runs without \code{dpkg}, without root privileges, without container
runtimes, and without consulting host system libraries.

\section{From Verilog to Cells}
\label{sec:t-cells}

Listing~\ref{lst:t-synth} shows the complete Yosys synthesis script. The
build interpolates the template placeholders delimited by \code{@},
adhering to the command sequence codified by the OpenROAD flow scripts.

\lstinputlisting[style=script, caption={Vreteno onto standard
cells.}, label={lst:t-synth}]{cpu/vreteno/asic/synth.ys}

Six of these commands perform critical structural transformations.
First, \code{dfflibmap} must execute before \code{abc}: sequential state
elements must be bound to library flip-flops prior to combinational
optimization, as the technology mapper cannot infer sequential cells.
Next, \code{setundef} assigns explicit binary levels to uninitialized or
don't-care nets, and \code{splitnets} decomposes multibit buses into
individual scalar wires—an essential flattening step, as gate-level
standard-cell netlists admit no vector abstractions. Then \code{hilomap}
substitutes literal logic constants (\code{1'b0} and \code{1'b1}) with
explicit tie-off cells (\code{TIEL} and \code{TIEH}); the mapper will
not select these autonomously because the cell library flags them with
the \code{dont\_use} attribute. Finally, \code{insbuf} inserts active
driver buffers on feedthrough ports connecting an input directly to an
output; in a physical layout, every driven wire requires an active cell
driver.

The ordering of these passes is deceptively unforgiving. Moving buffer
insertion ahead of bus splitting caused Yosys to instantiate buffers
across every constituent wire of every bus: 3,800 redundant cells that
served no electrical purpose, bloating the design by twenty percent.
This bloat generated no tool warnings; it was uncovered entirely
through manual post-synthesis cell accounting.

The resulting structural netlist comprises \asiccells{} standard cells
interconnected across \asicnets{} nets, of which \asicflops{} are
sequential flip-flops.

\subsection{The Memories}

This inventory contrasts sharply with Vreteno's footprint on the
Artix-7 FPGA: 2,151 lookup tables, 489 flip-flops, two block RAM tiles,
and four DSP48 arithmetic slices~\cite{vreteno}. The threefold expansion
in flip-flop count stems entirely from the register file. On the FPGA,
the 32-entry dual-read, single-write register file is mapped onto twelve
distributed RAM (LUTRAM) primitives, exploiting specialized fabric
features engineered for shallow multi-ported memories. On standard cells
without an SRAM macro, thirty-two 32-bit registers require 1,024
discrete D-type flip-flops. Combined with pipeline registers, CSRs, and
control logic, the ASIC count reaches \asicflops{} flip-flops, whereas
the board required only 489 because its register file was not built from
flip-flops at all.

The instruction memory underwent an inverse transformation. On the FPGA,
it occupies two 32-bit $\times$ 1024-word block RAM tiles, providing
concurrent dual-word fetches so that 16-bit compressed instructions
spanning a 32-bit word boundary can be decoded in a single cycle. On the
ASIC, lacking block RAM and an SRAM compiler, the read-only instruction
store resolves to pure combinational logic: because the program is static
and known at elaboration, synthesis collapses the memory into a massive
mask ROM. Its gate count is an exact function of binary entropy—a
different program produces a different gate count and a different
physical area. This distinction is vital to state explicitly: the
silicon area reported below measures this core executing this exact
benchmark binary, rather than a generic programmable microarchitecture.

The multiplier exemplifies a third trade-off. The four dedicated DSP48
slices utilized on the FPGA are replaced by an array multiplier
synthesized entirely from discrete standard cells.

\section{From Cells to a Layout}
\label{sec:t-pnr}

Listing~\ref{lst:t-sdc} specifies the timing constraints applied during
place-and-route. We constrain the clock period to 3\,ns (333\,MHz)—over
three times the FPGA board target (100\,MHz)—reflecting the inherently
faster intrinsic delays of 45\,nm standard cells over 28\,nm FPGA lookup
tables. Specifying an aggressive target gives meaning to downstream
timing slack: every picosecond of margin reflects real design
headroom. Primary I/O ports are timed assuming an external chip of
identical performance: input arrival delays and output load budgets are
each allocated 20\% of the clock period ($0.6$\,ns). Leaving boundary
ports unconstrained is the conventional recipe for timing reports that
are deceptively clean and practically meaningless.

\lstinputlisting[style=script, caption={What the core is asked to
meet.}, label={lst:t-sdc}]{cpu/vreteno/asic/vreteno.sdc}

The physical layout flow, executed by \code{cpu/vreteno/asic/pnr.tcl},
progresses in six sequential stages. First, floorplanning establishes a
square die boundary with standard-cell rows and routing tracks, placing
external I/O pins along the perimeter on upper metal layers clear of the
core. Next, well-tap cells are inserted every 120\,$\mu$m along cell rows
and at row boundaries to prevent latch-up, followed by power-grid
synthesis (PDN) providing horizontal standard-cell supply rails and
orthogonal straps on higher metal layers. Global placement then spreads
cells across the core to minimize wire length, followed by timing-driven
cell resizing to repair transition-time and fanout violations, and
legalization into continuous rows. Clock tree synthesis (CTS) constructs
a balanced buffer tree delivering clock edges to all \asicflops{}
flip-flops with minimal skew, after which localized placement refinement
shortens clock nets. Global and detailed routing assign signals to
metal tracks and vias while satisfying geometric design rules, followed
by post-route timing repair to heal delays introduced by real wire
parasitics. Last of all, non-functional filler cells are packed into
remaining row gaps to ensure substrate and well continuity across the
core.

\begin{figure*}[t]
\centering
\input{asicmap.tex}
\caption{The placed core, represented as a grid of \asicbins{} bins per
side across the die. On the left, cell density across all standard
cells prior to filler insertion: the pale border marks the clearance
margin between core rows and the die boundary. On the right, sequential
flip-flop density alone, normalized against the peak bin
(\asicffpeak\,\% of bin area); flip-flops comprise roughly a tenth of
total cells, so against a fully occupied square the panel would appear
as a uniform pale wash. The architectural state exhibits distinct
spatial geometry: regions devoid of flip-flops consist entirely of
combinational logic—dominated by the mask ROM encoding the program—while
dense flip-flop clusters reveal the register file and pipeline stage
registers.}
\label{fig:t-map}
\end{figure*}

\section{What the Run Says}
\label{sec:t-results}

Table~\ref{tab:t-results} summarizes the final physical design metrics.
Every datum in the table was emitted directly by OpenROAD into
machine-readable files and referenced in the LaTeX source via generated
macros, ensuring the narrative remains strictly synchronized with tool
outputs.

\begin{table}[htbp]
\caption{Vreteno on Nangate45, after routing.}
\label{tab:t-results}
\centering
\begin{tabular}{@{}lr@{}}
\toprule
Cells, after mapping & \asiccells{} \\
Nets & \asicnets{} \\
Flip-flops & \asicflops{} \\
Instances placed, before filler & \asicinstances{} \\
\addlinespace
Die & \asicdie{} $\times$ \asicdie\,$\mu$m \\
Cell area & \asicarea\,$\mu$m$^2$ \\
Core filled by cells & \asicutil\,\% \\
\addlinespace
Clock & \asicperiod\,ns (\asicfreq\,MHz) \\
Worst setup slack & \asicslack\,ns \\
Worst hold slack & \asicholdslack\,ns \\
Total negative slack & \asictns\,ns \\
\addlinespace
Routing violations & \asicdrc{} \\
\bottomrule
\end{tabular}
\end{table}

Timing closure is achieved with a tight setup margin: \asicslack\,ns of
slack on a \asicperiod\,ns clock cycle. Total Negative Slack (TNS) is
\asictns\,ns, demonstrating that the design contains no timing
violations across any path. Hold timing is likewise satisfied across all
register-to-register transfers. This distinction is critical: a setup
timing violation produces a chip that must run at reduced frequency,
whereas a hold violation produces silicon that fails catastrophically at
any operating speed.

Detailed routing concludes with \asicdrc{} design rule violations. In
physical verification, zero is the only acceptable outcome: an ASIC with
unrouted nets or short circuits is an unmanufacturable artifact.

Core utilization settles at \asicutil\,\%, an intentionally modest
figure. The initial floorplan dimensions were computed directly from
pre-placement cell area, and subsequent timing optimization inserted
repeaters and upsized gates that grew total area by roughly 25\%. While a
production ASIC pipeline would iteratively compress the core boundary to
minimize silicon cost, this reference flow prioritizes robust, zero-DRC
routing over aggressive silicon utilization.

\section{Where This Sits in the Build}
\label{sec:t-build}

Two Bazel targets govern this flow. The first,
\code{//cpu/vreteno/asic:vreteno\_cells}, executes logic synthesis and
technology mapping with Yosys. It completes in roughly ten seconds and is
included in the default build suite; any commit that introduces an HDL
construct incompatible with standard cells breaks the build immediately
rather than surfacing months later during tapeout.

The second target, \code{//cpu/vreteno/asic:vreteno\_pnr}, executes the
full OpenROAD place-and-route flow. Because detailed routing requires
twenty minutes of compute, the target is tagged \code{manual} (consistent
with our Vivado FPGA bitstream targets), keeping routine \code{bazel
test //...} runs fast enough for pre-commit verification.

Decoupling the 20-minute physical layout run from document compilation
poses a provenance challenge: how can a fast document build quote exact
layout metrics without manual transcription? The problem is identical to
tracking project timeline data, and it is resolved identically. The
physical run produces two compact text artifacts: the scalar timing
metrics and a spatial cell-density histogram. Both are checked into
source control. At document build time, \code{//tools/denmap} compiles
these text tables into LaTeX macros and generates the TikZ density plot
in Fig.~\ref{fig:t-map}. The final 40\,MB routed DEF remains an
uncommitted artifact in the build cache. As a result, no figure is a
bitmap pasted by hand, no quantitative claim can diverge from EDA
output, and updating any statistic requires executing the physical run
and checking in its emitted metrics.

\section{What a Tapeout Would Still Want}
\label{sec:t-not}

This flow terminates with a clean, routed design that closes timing
under nominal conditions. A genuine silicon tapeout begins there.

\emph{Foundry-certified physical verification.} Nangate45 provides
neither certified design rule decks nor calibrated device models. The
zero-DRC metric achieved here reflects the router's internal geometric
checks rather than independent signoff. Full physical verification
requires DRC and Layout-Versus-Schematic (LVS) validation using
foundry-supplied decks (such as KLayout or Calibre). On an open foundry
PDK with fabrication backing—such as SkyWater 130\,nm or
GlobalFoundries GF180MCU—these decks exist and would run directly
against the emitted layout.

\emph{A writable instruction store.} Implemented as a mask ROM,
Vreteno's program is frozen into silicon at fabrication. A general-purpose
microprocessor demands programmable execution memory, requiring either
an SRAM macro generated by a memory compiler or a boot ROM paired with a
serial interface (UART or SPI) to load code into memory before releasing
core reset. Vreteno already incorporates the bus interface needed to
attach such a loader.

\emph{I/O pad ring and ESD protection.} The core's \asiccells{} cells
currently terminate at perimeter pins on the core boundary—suitable for
macro integration in a larger SoC, but insufficient for a standalone
packaged die. A fabricated chip requires a pad ring: I/O pads with
level shifters, high-current output drivers, electrostatic discharge
(ESD) protection networks, bond or bump arrays, and independent I/O
supply rings.

\emph{Design for testability (DFT) and scan insertion.} The design
incorporates no scan chains; every sequential element is an ordinary
flip-flop. A manufactured part could be clocked functionally but not
shifted through scan vectors, rendering internal manufacturing defects
impossible to isolate. Inserting scan flip-flops and test-enable pins
is a standard synthesis transformation, representing the dividing line
between silicon that happens to run and silicon whose integrity can be
proven through automatic test pattern generation (ATPG).

\emph{Explicit power-on reset architecture.} Vreteno explicitly resets
its pipeline registers, CSRs, performance counters, and debug registers,
leaving its 32-entry register file uninitialized. On an FPGA, the
global configuration bitstream clears all storage elements at startup.
On silicon, unreset flip-flops power up in random states. While RISC-V
architectural semantics tolerate undefined values in registers
$x1$--$x31$ prior to software initialization, register $x0$ must
strictly evaluate to zero. Vreteno safeguards $x0$ by hardwiring reads
to zero rather than fetching from the storage array. Nonetheless,
physical signoff mandates explicit power-on-reset sequencing: an initial
value in Verilog is a simulation convenience, not a physical circuit.

\emph{Multi-corner, multi-mode (MCMM) timing signoff.} This run evaluated
a single operational point: typical process, nominal voltage, and
ambient room temperature (TT/1.1\,V/25\,$^\circ$C). Production signoff
requires analyzing PVT corners across the operating envelope: the slow
corner (SS, low voltage, high temperature) to guarantee setup timing at
target frequency, and the fast corner (FF, high voltage, sub-zero
temperature) to prevent catastrophic hold-time violations.

\emph{Parasitic extraction (PEX).} Wire delays in this run were modeled
using OpenROAD's internal resistance-capacitance estimations. Final
signoff demands rigorous 2.5D or 3D parasitic extraction from the
finished layout to generate a standard parasitic exchange format (SPEF)
file, followed by static timing analysis against extracted parasitics.

Each of these requirements is a known engineering task rather than an
architectural uncertainty. The physical flow demonstrated here
establishes the automated foundation onto which those signoff steps
attach.


\begin{thebibliography}{9}

\bibitem{vreteno}
F.~Filmar and D.~Janković, ``Vreteno: an RV32IMAC core in TxHDL,''
project repository \code{HDL/txhdl}, \code{//docs:vreteno}, 2026.

\bibitem{cover}
F.~Filmar and D.~Janković, ``TxHDL documents,'' project repository
\code{HDL/txhdl}, \code{//docs:cover}, 2026.

\end{thebibliography}

\end{document}
